\chapter{Cold Agglutinin Test}

\section*{What Is This Test?}
Cold-agglutinin testing detects antibodies that bind and agglutinate red blood cells at low temperatures. These antibodies can cause complement-mediated hemolysis and cold agglutinin disease.

\section*{Alternative Names}
Cold Agglutinin Titer; Cold Agglutinin Screen; Cold Autoantibody Test.

\section*{Why This Test Is Ordered}
Testing is used in suspected cold autoimmune hemolytic anemia, hemolysis after Mycoplasma or other infections, lymphoproliferative disorders, and unexplained red-cell agglutination or interference with blood grouping.

\section*{How This Test Is Performed}
Blood is collected and kept warm until serum or plasma is separated. The laboratory tests serial dilutions at defined temperatures and may determine the thermal amplitude—the highest temperature at which the antibody reacts. A direct antiglobulin test is often performed in parallel.

\section*{Preparation, Risks, and Limitations}
No fasting is required. Venipuncture is the usual risk. Improper cooling or warming during transport can invalidate the result. A high titer with a low thermal amplitude may be less clinically significant than a lower titer that reacts near body temperature.

\section*{Understanding Results}
A significant cold agglutinin, hemolysis, and complement-positive DAT support cold agglutinin disease, but infection, lymphoma, and other causes must be assessed. A positive test without hemolysis does not establish clinically important disease.

\section*{References}
British Society for Haematology guidance on autoimmune hemolytic anemia.

\clearpage
\section*{Understanding Results and Follow-up}
The laboratory report should identify the specimen, method, units, reference interval or decision limit, and whether the result is positive, negative, abnormal, or inconclusive. Reference intervals vary with age, sex, pregnancy, specimen type, assay, and laboratory; a result outside a range is not automatically a diagnosis, and a result within a range does not always exclude disease. Discuss unexpected results with the ordering clinician, who can decide whether repeat or confirmatory testing, treatment, or referral is appropriate. Do not change medicines or delay urgent care based on an isolated result.


\section*{Regulatory Status and Authorization Date}
Regulatory status applies to the specific assay, instrument, manufacturer, and intended use rather than to the test name alone. No single approval date can be assigned to this test category. For a commercial assay, verify the current product in the relevant regulator's database: the U.S. Food and Drug Administration (FDA) clearance, approval, or authorization; India's Central Drugs Standard Control Organisation (CDSCO) licence or permission; the European Union In Vitro Diagnostic Regulation (EU IVDR) CE certificate issued through the applicable conformity-assessment route; or the United Kingdom Medicines and Healthcare products Regulatory Agency (MHRA) registration and UKCA/CE status. Laboratory-developed tests may not have an individual FDA, CDSCO, EU IVDR, or MHRA approval date.
\clearpage
